\documentclass[12pt,a4paper]{exam}
\usepackage[utf8]{inputenc}
\usepackage{amsmath,amssymb,amsthm}
\usepackage[margin=1in]{geometry}
\usepackage{graphicx}
\usepackage[colorlinks=true]{hyperref}
\pagestyle{headandfoot}
\firstpageheader{MIT OpenCourseWare}{}{\textbf{EXTREMELY HARD -- MIT LEVEL}}
\runningheader{MIT OpenCourseWare}{}{\textbf{EXTREMELY HARD -- MIT LEVEL}}
\firstpagefooter{}{Page \thepage}{}
\runningfooter{}{Page \thepage}{}

\begin{document}

\begin{center}
{\LARGE \textbf{MIT OpenCourseWare}}\\21W.745\\Advanced Essay Workshop
\vspace{0.5cm}
{21W.745} --- {Advanced Essay Workshop}
\vspace{0.3cm}
May 2026
\vspace{0.3cm}
\textbf{EXTREMELY HARD -- MIT LEVEL}
\end{center}

\begin{center}
\textbf{Time Limit: 3 hours}
\end{center}

\begin{center}
\textbf{Total Marks: 100}
\end{center}

\vspace{0.5cm}

\begin{questions}

\question[3]
Analyze the limitations of standard approaches to Argumentation when applied to edge cases in 21W.745. Construct explicit counterexamples showing where intuitive reasoning fails.

\vspace{0.3cm}

\question[2]
Analyze the limitations of standard approaches to Rhetorical Analysis when applied to edge cases in 21W.745. Construct explicit counterexamples showing where intuitive reasoning fails.

\vspace{0.3cm}

\question[2]
Construct an original proof-level problem involving Expository Writing for 21W.745 (Advanced Essay Workshop) and provide a complete solution. Your problem should be at the level of a graduate qualifying exam.

\vspace{0.3cm}

\question[6]
Prove the fundamental theorem concerning Stylistics in the context of 21W.745 (Advanced Essay Workshop). State the theorem precisely, provide a complete proof, and discuss three important corollaries.

\vspace{0.3cm}

\question[3]
Construct an original proof-level problem involving Creative Nonfiction for 21W.745 (Advanced Essay Workshop) and provide a complete solution. Your problem should be at the level of a graduate qualifying exam.

\vspace{0.3cm}

\question[2]
Prove the fundamental theorem concerning Argumentation in the context of 21W.745 (Advanced Essay Workshop). State the theorem precisely, provide a complete proof, and discuss three important corollaries.

\vspace{0.3cm}

\question[3]
Prove the fundamental theorem concerning Rhetorical Analysis in the context of 21W.745 (Advanced Essay Workshop). State the theorem precisely, provide a complete proof, and discuss three important corollaries.

\vspace{0.3cm}

\question[6]
Prove the fundamental theorem concerning Expository Writing in the context of 21W.745 (Advanced Essay Workshop). State the theorem precisely, provide a complete proof, and discuss three important corollaries.

\vspace{0.3cm}

\question[2]
Construct an original proof-level problem involving Stylistics for 21W.745 (Advanced Essay Workshop) and provide a complete solution. Your problem should be at the level of a graduate qualifying exam.

\vspace{0.3cm}

\question[3]
Construct an original proof-level problem involving Creative Nonfiction for 21W.745 (Advanced Essay Workshop) and provide a complete solution. Your problem should be at the level of a graduate qualifying exam.

\vspace{0.3cm}

\question[5]
Prove the fundamental theorem concerning Argumentation in the context of 21W.745 (Advanced Essay Workshop). State the theorem precisely, provide a complete proof, and discuss three important corollaries.

\vspace{0.3cm}

\question[5]
Construct an original proof-level problem involving Rhetorical Analysis for 21W.745 (Advanced Essay Workshop) and provide a complete solution. Your problem should be at the level of a graduate qualifying exam.

\vspace{0.3cm}

\question[3]
Construct an original proof-level problem involving Expository Writing for 21W.745 (Advanced Essay Workshop) and provide a complete solution. Your problem should be at the level of a graduate qualifying exam.

\vspace{0.3cm}

\question[5]
Construct an original proof-level problem involving Stylistics for 21W.745 (Advanced Essay Workshop) and provide a complete solution. Your problem should be at the level of a graduate qualifying exam.

\vspace{0.3cm}

\question[5]
Construct an original proof-level problem involving Creative Nonfiction for 21W.745 (Advanced Essay Workshop) and provide a complete solution. Your problem should be at the level of a graduate qualifying exam.

\vspace{0.3cm}

\question[8]
Analyze the limitations of standard approaches to Argumentation when applied to edge cases in 21W.745. Construct explicit counterexamples showing where intuitive reasoning fails.

\vspace{0.3cm}

\question[4]
Prove the fundamental theorem concerning Rhetorical Analysis in the context of 21W.745 (Advanced Essay Workshop). State the theorem precisely, provide a complete proof, and discuss three important corollaries.

\vspace{0.3cm}

\question[6]
Construct an original proof-level problem involving Expository Writing for 21W.745 (Advanced Essay Workshop) and provide a complete solution. Your problem should be at the level of a graduate qualifying exam.

\vspace{0.3cm}

\question[5]
Construct an original proof-level problem involving Stylistics for 21W.745 (Advanced Essay Workshop) and provide a complete solution. Your problem should be at the level of a graduate qualifying exam.

\vspace{0.3cm}

\question[4]
Construct an original proof-level problem involving Creative Nonfiction for 21W.745 (Advanced Essay Workshop) and provide a complete solution. Your problem should be at the level of a graduate qualifying exam.

\vspace{0.3cm}

\question[3]
Analyze the limitations of standard approaches to Argumentation when applied to edge cases in 21W.745. Construct explicit counterexamples showing where intuitive reasoning fails.

\vspace{0.3cm}

\question[2]
Analyze the limitations of standard approaches to Rhetorical Analysis when applied to edge cases in 21W.745. Construct explicit counterexamples showing where intuitive reasoning fails.

\vspace{0.3cm}

\question[6]
Construct an original proof-level problem involving Expository Writing for 21W.745 (Advanced Essay Workshop) and provide a complete solution. Your problem should be at the level of a graduate qualifying exam.

\vspace{0.3cm}

\question[4]
Prove the fundamental theorem concerning Stylistics in the context of 21W.745 (Advanced Essay Workshop). State the theorem precisely, provide a complete proof, and discuss three important corollaries.

\vspace{0.3cm}

\question[3]
Analyze the limitations of standard approaches to Creative Nonfiction when applied to edge cases in 21W.745. Construct explicit counterexamples showing where intuitive reasoning fails.

\vspace{0.3cm}

\end{questions}

\newpage
\section*{Solutions}

\textbf{Solution 1:}

The edge case analysis reveals the subtle assumptions underlying standard treatments of Argumentation. By constructing explicit counterexamples, we see that the usual hypotheses cannot be weakened without loss of the theorem's conclusion.

\vspace{0.5cm}

\textbf{Solution 2:}

The edge case analysis reveals the subtle assumptions underlying standard treatments of Rhetorical Analysis. By constructing explicit counterexamples, we see that the usual hypotheses cannot be weakened without loss of the theorem's conclusion.

\vspace{0.5cm}

\textbf{Solution 3:}

Solution approach: First recall the definitions and key theorems related to Expository Writing from 21W.745. Then apply the appropriate techniques to construct a rigorous argument. The solution must be self-contained and reference only the material from the course.

\vspace{0.5cm}

\textbf{Solution 4:}

The proof follows from the definitions and properties established in 21W.745. The key insight is to apply the relevant theorem and verify the hypotheses hold. Detailed solution depends on the specific formulation of the theorem.

\vspace{0.5cm}

\textbf{Solution 5:}

Solution approach: First recall the definitions and key theorems related to Creative Nonfiction from 21W.745. Then apply the appropriate techniques to construct a rigorous argument. The solution must be self-contained and reference only the material from the course.

\vspace{0.5cm}

\textbf{Solution 6:}

The proof follows from the definitions and properties established in 21W.745. The key insight is to apply the relevant theorem and verify the hypotheses hold. Detailed solution depends on the specific formulation of the theorem.

\vspace{0.5cm}

\textbf{Solution 7:}

The proof follows from the definitions and properties established in 21W.745. The key insight is to apply the relevant theorem and verify the hypotheses hold. Detailed solution depends on the specific formulation of the theorem.

\vspace{0.5cm}

\textbf{Solution 8:}

The proof follows from the definitions and properties established in 21W.745. The key insight is to apply the relevant theorem and verify the hypotheses hold. Detailed solution depends on the specific formulation of the theorem.

\vspace{0.5cm}

\textbf{Solution 9:}

Solution approach: First recall the definitions and key theorems related to Stylistics from 21W.745. Then apply the appropriate techniques to construct a rigorous argument. The solution must be self-contained and reference only the material from the course.

\vspace{0.5cm}

\textbf{Solution 10:}

Solution approach: First recall the definitions and key theorems related to Creative Nonfiction from 21W.745. Then apply the appropriate techniques to construct a rigorous argument. The solution must be self-contained and reference only the material from the course.

\vspace{0.5cm}

\textbf{Solution 11:}

The proof follows from the definitions and properties established in 21W.745. The key insight is to apply the relevant theorem and verify the hypotheses hold. Detailed solution depends on the specific formulation of the theorem.

\vspace{0.5cm}

\textbf{Solution 12:}

Solution approach: First recall the definitions and key theorems related to Rhetorical Analysis from 21W.745. Then apply the appropriate techniques to construct a rigorous argument. The solution must be self-contained and reference only the material from the course.

\vspace{0.5cm}

\textbf{Solution 13:}

Solution approach: First recall the definitions and key theorems related to Expository Writing from 21W.745. Then apply the appropriate techniques to construct a rigorous argument. The solution must be self-contained and reference only the material from the course.

\vspace{0.5cm}

\textbf{Solution 14:}

Solution approach: First recall the definitions and key theorems related to Stylistics from 21W.745. Then apply the appropriate techniques to construct a rigorous argument. The solution must be self-contained and reference only the material from the course.

\vspace{0.5cm}

\textbf{Solution 15:}

Solution approach: First recall the definitions and key theorems related to Creative Nonfiction from 21W.745. Then apply the appropriate techniques to construct a rigorous argument. The solution must be self-contained and reference only the material from the course.

\vspace{0.5cm}

\textbf{Solution 16:}

The edge case analysis reveals the subtle assumptions underlying standard treatments of Argumentation. By constructing explicit counterexamples, we see that the usual hypotheses cannot be weakened without loss of the theorem's conclusion.

\vspace{0.5cm}

\textbf{Solution 17:}

The proof follows from the definitions and properties established in 21W.745. The key insight is to apply the relevant theorem and verify the hypotheses hold. Detailed solution depends on the specific formulation of the theorem.

\vspace{0.5cm}

\textbf{Solution 18:}

Solution approach: First recall the definitions and key theorems related to Expository Writing from 21W.745. Then apply the appropriate techniques to construct a rigorous argument. The solution must be self-contained and reference only the material from the course.

\vspace{0.5cm}

\textbf{Solution 19:}

Solution approach: First recall the definitions and key theorems related to Stylistics from 21W.745. Then apply the appropriate techniques to construct a rigorous argument. The solution must be self-contained and reference only the material from the course.

\vspace{0.5cm}

\textbf{Solution 20:}

Solution approach: First recall the definitions and key theorems related to Creative Nonfiction from 21W.745. Then apply the appropriate techniques to construct a rigorous argument. The solution must be self-contained and reference only the material from the course.

\vspace{0.5cm}

\textbf{Solution 21:}

The edge case analysis reveals the subtle assumptions underlying standard treatments of Argumentation. By constructing explicit counterexamples, we see that the usual hypotheses cannot be weakened without loss of the theorem's conclusion.

\vspace{0.5cm}

\textbf{Solution 22:}

The edge case analysis reveals the subtle assumptions underlying standard treatments of Rhetorical Analysis. By constructing explicit counterexamples, we see that the usual hypotheses cannot be weakened without loss of the theorem's conclusion.

\vspace{0.5cm}

\textbf{Solution 23:}

Solution approach: First recall the definitions and key theorems related to Expository Writing from 21W.745. Then apply the appropriate techniques to construct a rigorous argument. The solution must be self-contained and reference only the material from the course.

\vspace{0.5cm}

\textbf{Solution 24:}

The proof follows from the definitions and properties established in 21W.745. The key insight is to apply the relevant theorem and verify the hypotheses hold. Detailed solution depends on the specific formulation of the theorem.

\vspace{0.5cm}

\textbf{Solution 25:}

The edge case analysis reveals the subtle assumptions underlying standard treatments of Creative Nonfiction. By constructing explicit counterexamples, we see that the usual hypotheses cannot be weakened without loss of the theorem's conclusion.

\vspace{0.5cm}

\end{document}